\section*{Capítulo \ref{ch:b2:plant-flowering} --- Desarrollo reproductor y control de la floración}

\begin{solution}{exo:b2:plant-flowering:1}
\omterm{def:b2:plant-meristems:meristem}{Meristemo} vegetativo: forma hojas con yemas axilares, indefinidamente
(crecimiento indefinido). \omterm{def:b2:plant-flowering:transition}{Meristemo de inflorescencia}: forma brácteas y
\omterm{def:b2:plant-meristems:meristem}{meristemos} florales en sus axilas; de crecimiento indefinido en muchas
especies. \omterm{def:b2:plant-meristems:meristem}{Meristemo} floral: forma los cuatro verticilos y se detiene ---
de crecimiento definido.
\end{solution}

\begin{solution}{exo:b2:plant-flowering:2}
Las \omterm{prop:b2:plant-flowering:photoperiod}{plantas de día corto} florecen cuando el día es más corto que una
duración crítica (crisantemo, soja, arroz); las de día largo, cuando es
más largo (espinaca, trigo, \emph{Arabidopsis}); las neutras florecen al
alcanzar un tamaño (tomate, maíz). Una \omterm{prop:b2:plant-flowering:photoperiod}{planta de día corto} mide la
duración de la noche ininterrumpida.
\end{solution}

\begin{solution}{exo:b2:plant-flowering:3}
Tipo silvestre: sépalo, pétalo, \omterm{def:b2:angiosperm-reproduction:flower}{estambre}, \omterm{def:b2:angiosperm-reproduction:flower}{carpelo}. Sin A: \omterm{def:b2:angiosperm-reproduction:flower}{carpelo},
\omterm{def:b2:angiosperm-reproduction:flower}{estambre}, \omterm{def:b2:angiosperm-reproduction:flower}{estambre}, \omterm{def:b2:angiosperm-reproduction:flower}{carpelo}. Sin B: sépalo, sépalo, \omterm{def:b2:angiosperm-reproduction:flower}{carpelo}, \omterm{def:b2:angiosperm-reproduction:flower}{carpelo}. Sin
C: sépalo, pétalo, pétalo, sépalo, y la \omterm{def:b2:angiosperm-reproduction:flower}{flor} se repite en su interior.
\end{solution}

\begin{solution}{exo:b2:plant-flowering:4}
Semanas de frío y humedad (el invierno ha pasado); la luz (la \omterm{def:b2:angiosperm-reproduction:seed}{semilla} está
cerca de la superficie); las temperaturas alternas (suelo desnudo, sin
cubierta vegetal); el lavado por la lluvia (han llegado las lluvias); la
escarificación (paso por un tubo digestivo o abrasión en el suelo); el
humo y el calor (un incendio ha despejado el terreno).
\end{solution}

\begin{solution}{exo:b2:plant-flowering:5}
16/8: noche de 8 horas, por debajo de las 10 críticas: no. 12/12: sí.
12 de luz, 6 de oscuridad, destello, 6 de oscuridad: dos noches de 6
horas: no. 8 de luz, 4 de oscuridad, 4 de luz, 8 de oscuridad: noche más
larga de 8 horas: no.
\end{solution}

\begin{solution}{exo:b2:plant-flowering:6}
El \omterm{def:b2:plant-flowering:florigen}{florígeno} es una señal transmisible por injerto, fabricada en la hoja e
igual en las especies de día corto y de día largo --- la diferencia está
en cuándo la fabrica la hoja, no en qué fabrica. Con el floema
interrumpido, el injerto no induce la floración: la señal viaja por el
floema.
\end{solution}

\begin{solution}{exo:b2:plant-flowering:7}
$n = \ln(1/0.15)/0.05 = 1.90/0.05 = 38$ días fríos. Los periodos suman
$47$: la planta queda vernalizada el día 11 del tercer periodo. Con
$\alpha = 0.02$ necesitaría 95 días y no lo quedaría.
\end{solution}

\begin{solution}{exo:b2:plant-flowering:8}
El gen del represor queda silenciado por marcas de la cromatina que se
copian en cada mitosis, de modo que todas las células del tallo recuerdan
el invierno; en la línea germinal y en el embrión temprano las marcas se
borran y el gen vuelve a expresarse. Es útil porque una \omterm{def:b2:angiosperm-reproduction:seed}{semilla} liberada en
verano no debe heredar el permiso de floración de su progenitora: debe
esperar a un invierno propio.
\end{solution}

\begin{solution}{exo:b2:plant-flowering:9}
R: germinan (el fitocromo pasa a la forma que absorbe el rojo lejano,
Pfr, la activa). R y después RL: no (vuelve a Pr). R, RL, R: germinan. R,
RL, R, RL: no. Solo cuenta el último destello, porque cada conversión es
completa y la \omterm{def:b2:angiosperm-reproduction:seed}{semilla} lee el estado final del pigmento.
\end{solution}

\begin{solution}{exo:b2:plant-flowering:10}
$p = 0.7$: $r = 1.28$, $1.42$, $1.40$, $-\infty$ para $g = 0.4, 0.6, 0.8,
1$: el óptimo, cerca de $g = 0.7$. $p = 0.95$: $1.99$, $2.33$, $2.55$,
$-\infty$: el óptimo se desplaza hacia $g = 1$ (hasta cerca de $0.95$), con
una reserva solo simbólica. Los desiertos, donde los buenos años son
raros, favorecen una \omterm{def:b2:angiosperm-reproduction:seed}{latencia} intensa y un gran banco de \omterm{def:b2:angiosperm-reproduction:seed}{semillas}; las
praderas, donde la mayoría de los años son buenos, favorecen germinar casi
todo.
\end{solution}

\begin{solution}{exo:b2:plant-flowering:11}
B en todos los verticilos: pétalo, pétalo, \omterm{def:b2:angiosperm-reproduction:flower}{estambre}, \omterm{def:b2:angiosperm-reproduction:flower}{estambre}. C en todos
los verticilos (C reprime A): \omterm{def:b2:angiosperm-reproduction:flower}{carpelo}, \omterm{def:b2:angiosperm-reproduction:flower}{estambre}, \omterm{def:b2:angiosperm-reproduction:flower}{estambre}, \omterm{def:b2:angiosperm-reproduction:flower}{carpelo}. Sin A
ni B: C sola en todas partes --- \omterm{def:b2:angiosperm-reproduction:flower}{carpelos} en los cuatro verticilos. C
también detiene el \omterm{def:b2:plant-meristems:meristem}{meristemo}; sin ella, el centro sigue siendo de
crecimiento indefinido y forma una \omterm{def:b2:angiosperm-reproduction:flower}{flor} dentro de otra.
\end{solution}

\begin{solution}{exo:b2:plant-flowering:12}
La duración del día se integra de una noche a otra mediante la
coincidencia de la luz con el reloj, y la inducción requiere varias de
esas noches; el frío se integra en forma de marcas de cromatina que se
acumulan; el banco de \omterm{def:b2:angiosperm-reproduction:seed}{semillas} integra a lo largo de los años
escalonando la germinación. Un integrador ignora un único día anómalo, un
deshielo de enero o un mal año aislado; un umbral sobre el valor
instantáneo haría florecer la planta en una semana cálida de febrero, o
germinar todas las \omterm{def:b2:angiosperm-reproduction:seed}{semillas} con la primera lluvia.
\end{solution}

\begin{solution}{pb:b2:plant-flowering:1}
\textbf{1.} Una noche de 8 horas es más corta que las 9,5 horas
críticas: no hay floración. Para vender en agosto, el cultivador debe
alargar artificialmente las noches con una tela negra.
\textbf{2.} Una noche de 12 horas. Una apertura el 15 de agosto supone una
inducción completada 8 semanas antes, hacia el 20 de junio; las doce
noches inductivas deben empezar hacia el 8 de junio.
\textbf{3.} Iluminar el cultivo unos minutos en mitad de cada noche: una
noche de 14,5 horas se convierte en dos de unas 7 horas, ambas por debajo
de la duración crítica. Basta un destello porque el fitocromo se
convierte en minutos y la planta solo mide la duración de la oscuridad
ininterrumpida.
\textbf{4.} Un destello a las 21:00 deja un periodo oscuro de las 21:00 a
las 07:30, 10,5 horas --- más que 9,5: ese tramo induce. La interrupción
debe situarse de modo que ningún tramo supere la duración crítica.
\textbf{5.} Una noche perdida rompe la serie de doce noches inductivas
consecutivas: el recuento vuelve a empezar y la floración se retrasa unos
días, tantos como se hubieran acumulado. Un cadillo necesita una sola
noche inductiva y ya habría quedado comprometido.
\textbf{6.} El fitocromo no absorbe la luz verde, de modo que la planta no
ve el destello; el rojo lo convierte a la forma activa, que se lee como
día; el rojo lejano inmediatamente después lo devuelve a la forma
inicial, de modo que la noche no queda interrumpida.
\textbf{7.} La \omterm{prop:b2:plant-flowering:photoperiod}{planta de día largo} florece: el \omterm{def:b2:plant-flowering:florigen}{florígeno} atraviesa la unión
del injerto por el floema. Esto identifica una señal móvil, transmisible
por injerto, común a ambas clases fotoperiódicas.
\textbf{8.} $n = \ln 10/0.06 = 38.4$: 39 días fríos.
\textbf{9.} 10 en noviembre, 35 al final de enero, y el día frío número 39
cae en el cuarto día frío de febrero: vernalizado a principios de
febrero.
\textbf{10.} Seis días: $f = e^{-0.36} = 0.70 > 0.10$: permanece vegetativo
todo el verano y no da grano --- el trigo de invierno debe sembrarse en
otoño.
\textbf{11.} Sin el represor florece sin frío: puede sembrarse en
primavera y cosecharse el mismo verano --- un trigo de primavera.
\textbf{12.} Las marcas de silenciamiento de la cromatina del represor se
copian en cada mitosis, de modo que el ápice lo recuerda a lo largo de la
primavera; en la \omterm{def:b2:meiosis-heredity:meiosis}{meiosis} y en el embrión las marcas se borran, y cada
generación debe contar su propio invierno.
\textbf{13.} La memoria debe residir en las células que formarán la \omterm{def:b2:angiosperm-reproduction:flower}{flor} y
que persisten a lo largo del invierno --- el \omterm{def:b2:plant-meristems:meristem}{meristemo} ---, y es un estado
de su cromatina, no una señal móvil; las hojas que miden la duración del
día se caen, y exportan en cambio una proteína.
\textbf{14.} Clase C: sépalo, pétalo, pétalo, sépalo (A se extiende al cuarto
verticilo y la \omterm{def:b2:angiosperm-reproduction:flower}{flor} se repite en su interior).
\textbf{15.} Sin \omterm{def:b2:angiosperm-reproduction:flower}{estambres} no hay \omterm{def:b2:plant-life-cycles:heterospory}{polen}; algunos \omterm{def:b2:angiosperm-reproduction:flower}{carpelos} se forman donde C
tiene una actividad débil, de modo que se forman \omterm{def:b2:angiosperm-reproduction:seed}{semillas} de vez en
cuando.
\textbf{16.} Clase A: \omterm{def:b2:angiosperm-reproduction:flower}{carpelo}, \omterm{def:b2:angiosperm-reproduction:flower}{estambre}, \omterm{def:b2:angiosperm-reproduction:flower}{estambre}, \omterm{def:b2:angiosperm-reproduction:flower}{carpelo}.
\textbf{17.} C reprime A: el verticilo 1 se convierte en \omterm{def:b2:angiosperm-reproduction:flower}{carpelo}, el
verticilo 2 (B y C) en \omterm{def:b2:angiosperm-reproduction:flower}{estambre} --- \omterm{def:b2:angiosperm-reproduction:flower}{carpelo}, \omterm{def:b2:angiosperm-reproduction:flower}{estambre}, \omterm{def:b2:angiosperm-reproduction:flower}{estambre}, \omterm{def:b2:angiosperm-reproduction:flower}{carpelo},
la \omterm{def:b2:angiosperm-reproduction:flower}{flor} del mutante A.
\textbf{18.} Los genes ABC son una familia de \omterm{prop:b2:cell-differentiation:transcriptional}{factores de transcripción}
heredada del antepasado común de todas las plantas con \omterm{def:b2:angiosperm-reproduction:flower}{flores}; el módulo
ya estaba establecido antes de que divergieran las rosas, las bocas de
dragón y \emph{Arabidopsis}, y cada una lo ha conservado.
\textbf{19.} B sin A ni C no especifica ningún órgano floral: un órgano de
tipo foliar o de tipo bráctea.
\textbf{20.} $p = 0.8$: $r = 1.44$ ($g = 0.5$), $1.59$ ($0.7$), $1.55$
($0.9$).
\textbf{21.} De los tres, $g = 0.7$ (el verdadero óptimo está cerca de
$0.8$); con $g = 1$ el factor de los malos años es cero y el linaje se
extingue con el primer mal año.
\textbf{22.} $p = 0.5$: $0.51$, $0.41$, $-0.03$: $g = 0.5$ es el mejor, y
un valor aún menor sería mejor --- más \omterm{def:b2:angiosperm-reproduction:seed}{latencia} en el desierto.
\textbf{23.} Las \omterm{def:b2:angiosperm-reproduction:seed}{semillas} del cultivador se almacenan, no se entierran, y
sobreviven con certeza; no necesita diversificar el riesgo dentro del
campo, pero debe guardar una reserva en el granero para volver a sembrar
tras un año fallido --- la misma reserva, en manos del cultivador en lugar
de en las del suelo.
\textbf{24.} Una \omterm{def:b2:angiosperm-reproduction:seed}{semilla} pequeña lleva reservas para un tallo de uno o dos
centímetros; si germinara en la oscuridad bajo tierra moriría antes de
alcanzar la luz, de modo que espera a la luz --- que solo le llega cerca de
la superficie. El arado sube a la superficie las \omterm{def:b2:angiosperm-reproduction:seed}{semillas} enterradas y
desencadena una oleada de malas hierbas.
\textbf{25.} Noche crítica de 9,5 horas, cubrición desde el 8 de junio
aproximadamente para abrir el 15 de agosto; 39 días fríos; mejor
$g \approx 0.8$ para $p = 0.8$ y alrededor de $0.5$ para $p = 0.5$.
\end{solution}
